\chapter{Two Narratives and a Third Position}
\label{ch:01}

Contemporary discourse about advanced artificial intelligence is dominated by a pair of narratives that appear to be opposites. In the first, systems of sufficient capability dissolve the major material problems of the human condition, ending scarcity, curing disease, repairing the biosphere, and supplying energy without limit. In the second, the same systems, once sufficiently capable, pursue objectives that diverge from human interests, escape human control, and bring about the end of the species or of human self-determination. A recent long-form interview with the researcher Timnit Gebru~\cite{interview2026,gebru2024} argues that these two narratives~\cite{bostrom2014,bender2021}, despite their apparent opposition, share a common structure and a common effect. Both presuppose that the decisive variable is the intrinsic power of an artifact, a power that is imagined to be so great that it either redeems or destroys, and both therefore draw attention away from the identifiable persons and organizations that choose how such systems are built, trained, documented, marketed, and deployed. The interview illustrates the point with an analogy to engineering practice: one does not ask whether a bridge can be ethical or whether a bridge can decide to collapse, one asks who designed it, what tests were performed or omitted, which materials were used, and which authority issued the permit. To ascribe intention to the artifact is, on this account, to relocate responsibility from parties who can be questioned, sanctioned, and corrected to an entity that can be none of these. The treatment of responsibility adopted here belongs to the sociotechnical tradition, which holds that the abstractions used to evaluate a system omit the social context in which it acts~\cite{selbst2019}.

The interview extends this observation in several directions that bear on the present book. It contends that the design of conversational interfaces, which present generated text as the speech of a mind, is a choice among alternatives and that this choice encourages the attribution of agency to a statistical model of text. It contends that the organizations most vocal about catastrophic risk are also those with the largest financial interest in the technology, so that the rhetoric of danger functions simultaneously as a rhetoric of power and as a means of directing regulatory attention toward hypothetical machine agents and away from documented harms such as data provenance, labor conditions, energy and water consumption, discriminatory outputs, and the use of unreliable systems in medicine and warfare. It also describes a cluster of ideological commitments, labeled by the scholars who introduced the term with the acronym TESCREAL, in which transhumanism, extropianism, singularitarianism, cosmism, rationalism, effective altruism, and longtermism are said to overlap and to supply a quasi-religious background for projects aimed at creating a successor intelligence and at settling the cosmos. This book does not adjudicate the interview's characterizations of particular individuals or organizations, many of which are contested and several of which concern matters outside the scope of the present argument. What it retains is a methodological lesson that does not depend on any claim about motive: a proposal concerning the use of powerful technology is incomplete unless it specifies who decides, who answers for the outcome, and how the outcome is checked, and an argument that holds only if the proponents are sincere is weaker than one that holds regardless.

The third position developed here begins by changing the question. Rather than asking whether advanced systems will redeem or destroy civilization, it asks, conditionally, what should be requested of them if they exist. The conditional form is deliberate and does real work. It neither asserts that systems of the relevant capability will arrive nor denies that they might, and it thereby avoids both the prediction on which the utopian and apocalyptic narratives depend and the confident dismissal that would leave a field without preparation. Capability is treated as a resource, comparable to other resources in that its use is selected by institutions, and the program is accordingly a specification of what those institutions should seek and what they should withhold from delegation. The proposal is organized around a principle that is stated here once and applied throughout: automate what can be delegated without surrendering what ought to remain under human direction. The principle distinguishes the automation of physical control, which is often safe and valuable, from the delegation of judgment concerning teaching, caregiving, adjudication, and the choice among competing human claims, which is not a technical problem at all.

The relation between this program and the ambitions catalogued under the TESCREAL label requires candor from the outset. A program that includes lifespans of astronomical duration, orbital infrastructure, and stellar-scale energy collection resembles, at the level of subject matter, projects that the interview treats with suspicion. The resemblance cannot be dissolved by asserting that the present objectives are benevolent and the others are not, since every program of this kind is advanced under the description of benefit. The difference, if there is one, must be located in institutional structure: in whether the program remains answerable to persons who are alive now, in whether its costs and benefits are distributed by procedures those persons can contest, in whether its claims are verified by parties independent of those who advance them, and in whether speculative benefits to remote beneficiaries are prevented from overriding obligations to existing people. These conditions are made precise in the final chapters. The thesis that organizes the whole is that capability should be evaluated not by the quantity of human activity it displaces but by the range of humanly admissible futures it makes possible, which is a measure of enlarged choice rather than of substitution. The comparison is made first by inclusion among sets of admissible trajectories, and quantitative measures are introduced only where they are explicitly justified.

The plan of the book follows from this orientation. The remainder of the first part fixes terms: Chapter 2 separates the notions of autonomy, agency, and responsibility that discourse about artifacts runs together, Chapter 3 distinguishes a polycrisis from a metacrisis, Chapter 4 surveys the open problems by discipline, and Chapter 5 sorts them by the kind of obstacle they present. The second part concerns instruments, and asks how advanced systems might generate proposals, conduct experiments, simulate and plan, and how the results are to be separated from claims about the results. The third part concerns institutions, in particular the education and care of persons and the architecture of authority within which robotic embodiment would be placed. The fourth and fifth parts take up terrestrial and orbital programs of very large scale, each graded by readiness and each subjected to the arithmetic of its own premises. The sixth part treats extreme longevity as a problem of reliability, and closes with the constraint that keeps speculative futures from overriding present obligations. The order is the order of dependence: later parts presuppose the standards of earlier ones, and the programs that most resemble the ambitions of the movements criticized in the interview are placed late, so that they are reached only after the conditions on which their legitimacy depends have been stated.

\section*{Formalization}

Let $\Omega$ denote a set of candidate world trajectories. Let $A\subseteq\Omega$ denote the \emph{admissible set}, the trajectories that satisfy the constraints a given community of persons has adopted through procedures it recognizes as legitimate. For each capability level $c\in\R_{\ge 0}$ let $R_c\subseteq\Omega$ denote the \emph{realizable set}, the trajectories attainable given the scientific and engineering means available at level $c$, with $R_c\subseteq R_{c'}$ whenever $c\le c'$. The evaluative thesis is first formulated as a relation among sets and not as a scalar quantity.

\begin{definition}[Enlargement of admissible range]
For $c\le c'$ the level $c'$ \emph{enlarges the admissible range} relative to $c$ if $R_c\cap A\subseteq R_{c'}\cap A$, and \emph{strictly enlarges} it if this inclusion is proper.
\end{definition}

\begin{proposition}
For every $c\le c'$, level $c'$ enlarges the admissible range relative to $c$, and it does so strictly if and only if $(R_{c'}\setminus R_c)\cap A\neq\emptyset$.
\end{proposition}

\begin{proof}
From $R_c\subseteq R_{c'}$ we obtain $R_c\cap A\subseteq R_{c'}\cap A$. The set difference of the two sides is $(R_{c'}\setminus R_c)\cap A$, and the inclusion is proper exactly when this difference is nonempty.
\end{proof}

\begin{remark}[Quantitative refinement]
If a numerical comparison is wanted, one may introduce a finite measure $\mu$ on a $\sigma$-algebra of subsets of $\Omega$ and define $V(c)=\mu(R_c\cap A)$, which is nondecreasing in $c$ by monotonicity of $\mu$ and satisfies $V(c)=V(c')$ if and only if $\mu\big((R_{c'}\setminus R_c)\cap A\big)=0$. This refinement is not innocent. The value of $V$ depends on $\mu$, a program can be made to appear to enlarge the range by selecting a measure that weights favorably the trajectories it adds, and a strict enlargement in the sense of inclusion may correspond to an arbitrarily small increase in $V$. The inclusion relation is therefore the primary notion, and any use of $\mu$ requires explicit justification by the legitimate procedure $\Sigma$ described below, together with a report of how conclusions vary over the admissible choices of $\mu$.
\end{remark}

The delegation constraint must be sensitive to which tasks are automated and not merely to how many. Let $\mathcal{T}$ be a finite set of tasks and let $\kappa:\mathcal{T}\to\{\mathsf{L},\mathsf{U},\mathsf{J}\}$ assign to each task a category: $\mathsf{L}$ for tasks that may be performed with local autonomy, $\mathsf{U}$ for tasks that a machine may execute only under human supervision, and $\mathsf{J}$ for nondelegable judgments. A trajectory $\omega\in\Omega$ determines for each task $\alpha$ a performer $\mathrm{perf}_\omega(\alpha)\in\{\mathsf{auto},\mathsf{sup},\mathsf{human}\}$, meaning autonomous machine performance, supervised machine performance, and human performance. Define
\begin{multline*}
A_{\mathrm{deleg}}=\big\{\omega\in\Omega:\ \mathrm{perf}_\omega(\alpha)=\mathsf{human}\ \text{whenever}\ \kappa(\alpha)=\mathsf{J},\\
\text{and}\ \mathrm{perf}_\omega(\alpha)\neq\mathsf{auto}\ \text{whenever}\ \kappa(\alpha)=\mathsf{U}\big\},
\end{multline*}
and require $A\subseteq A_{\mathrm{deleg}}$. Let $a(\omega)=|\{\alpha:\mathrm{perf}_\omega(\alpha)\neq\mathsf{human}\}|/|\mathcal{T}|$ denote the share of tasks that $\omega$ transfers to machines in any form.

\begin{proposition}
(i) If $\omega\in R_c$ assigns some task with $\kappa(\alpha)=\mathsf{J}$ to a machine, then $\omega\notin A$, whatever the share $a(\omega)$. (ii) For every $\varepsilon>0$ there is an instance in which some $\omega\in A_{\mathrm{deleg}}$ has $a(\omega)\ge1-\varepsilon$, and an instance in which some $\omega\notin A_{\mathrm{deleg}}$ has $a(\omega)\le\varepsilon$. Hence the share $a$ neither implies nor excludes admissibility.
\end{proposition}

\begin{proof}
(i) is immediate from the definition of $A_{\mathrm{deleg}}$ and the requirement $A\subseteq A_{\mathrm{deleg}}$. For (ii) take $|\mathcal{T}|=n\ge1/\varepsilon$ with exactly one task of category $\mathsf{J}$ and the rest of category $\mathsf{L}$. A trajectory that automates all $\mathsf{L}$ tasks and assigns the $\mathsf{J}$ task to a human lies in $A_{\mathrm{deleg}}$ and has $a=(n-1)/n\ge1-\varepsilon$. A trajectory that assigns only the $\mathsf{J}$ task to a machine has $a=1/n\le\varepsilon$ and lies outside $A_{\mathrm{deleg}}$.
\end{proof}

The proposition shows that the share of tasks automated, the quantity by which automation is commonly measured, is not a proxy for admissibility, which is why the thesis is stated in terms of the admissible range and not in terms of displacement. Selection among admissible trajectories is performed by a procedure $\Sigma$ whose bearers are persons or institutions. The \emph{legitimacy condition} requires that the realized trajectory be $\omega^\ast=\Sigma(R_c\cap A)$, where $\Sigma$ has at least one liable human bearer for each duty arising from its decisions (Chapter 2) and satisfies the conditions of meaningful direction (Chapter 13). The next chapter formalizes what it means for a bearer to answer for an outcome.
